\documentclass[11pt]{amsart}
\usepackage{amsmath,amssymb,amsthm,mathtools}
\usepackage[margin=1in]{geometry}
\usepackage{array}
\usepackage{hyperref}

\newtheorem{theorem}{Theorem}[section]
\newtheorem{proposition}[theorem]{Proposition}
\newtheorem{lemma}[theorem]{Lemma}
\newtheorem{corollary}[theorem]{Corollary}
\theoremstyle{definition}
\newtheorem{definition}[theorem]{Definition}
\newtheorem{criterion}[theorem]{Criterion}
\theoremstyle{remark}
\newtheorem{remark}[theorem]{Remark}

\DeclareMathOperator{\NS}{NS}
\DeclareMathOperator{\TP}{TP}
\newcommand{\bbZ}{\mathbb{Z}}
\newcommand{\bbQ}{\mathbb{Q}}
\newcommand{\bbF}{\mathbb{F}}
\newcommand{\bbR}{\mathbb{R}}
\newcommand{\Xib}{\Xi}
\newcommand{\LP}{\mathrm{LP}}
\newcommand{\RH}{\mathrm{RH}}

\title[Positivity formulations of RH and the function-field analogy]{Positivity
formulations of the Riemann hypothesis and the function-field analogy: a
heuristic checklist and some remarks}
\author{David Alejandro Trejo Pizzo}
\thanks{Independent Researcher. \texttt{dtrejopizzo@gmail.com}}


\begin{document}
\begin{abstract}
This is an expository note; it does not attempt to prove the Riemann
hypothesis (RH). We propose a four-condition heuristic checklist $D_0$ for
mechanisms that could exclude zeros off the critical line, and record how
eighteen known programmes fare against it; the assessment is a judgment, not a
theorem. In this assessment the only mechanisms meeting all four conditions
are the proofs of RH over finite fields, and the condition most often not met
is that the required bound be a corollary of a general theorem proved
independently for a whole class of objects, rather than a property of the
object at hand. We recall Schoenberg's total-positivity form of RH (RH holds
if and only if the Fourier transform of $1/\xi(\tfrac12+y)$ is a P\'olya
frequency function) and observe that P\'olya's kernel $\Phi$ itself is not a
P\'olya frequency function. We recall that in the Jensen-polynomial hierarchy
for $\Xi$ each fixed degree is known for large shift while the full hierarchy
is equivalent to RH, and we discuss the analogous pattern in other positivity
formulations. Finally we observe that in the function-field case the
positivity is a finite condition because the ordinates of the zeros are
periodic with period $2\pi/\log q$, and that no such periodicity is possible
for $\zeta$, since the number of distinct ordinates of zeros on the critical
line up to height $T$ is $\gg T\log T$. None of this proves or disproves RH.
\end{abstract}
\maketitle

\section{Introduction}
Many reformulations of the Riemann hypothesis and several programmes aimed at
it are known; general accounts are \cite{bgTitchmarsh,bgEdwards,bgConrey03,
bgBombieri00}. This note does not add a reformulation. It compares the
programmes with the one setting in which an analogue of RH is a theorem,
namely zeta functions over finite fields, and records three things: a
checklist (\S\ref{sec:disc}), the relation between conditions of finite order
and the full condition in positivity formulations of RH (\S\ref{sec:gap}), and
the role of the periodicity of the zeros over finite fields
(\S\ref{sec:fin}).

The checklist $D_0$ of \S\ref{sec:disc} is a heuristic; the mathematical
content is Theorem~\ref{thm:tp}, Proposition~\ref{prop:orders},
Proposition~\ref{prop:finit} and Proposition~\ref{thm:LI}. The first three are
classical, and the fourth is an immediate consequence of known results on the
density of zeros on the critical line. No statement below assumes RH, and none
implies it.

\section{The checklist \texorpdfstring{$D_0$}{D0}}\label{sec:disc}

Fix the completed zeta function
$\xi(s)=\tfrac12 s(s-1)\pi^{-s/2}\Gamma(s/2)\zeta(s)$, entire of order $1$,
with zeros $\tfrac12+i\gamma$; RH is the statement that every $\gamma$ is
real. A \emph{mechanism} is a mathematical structure $X$ together with an
asserted implication ``[a property of $X$] $\Rightarrow$ [no zeros off the
critical line, in some region].''

\begin{criterion}[$D_0$]\label{crit:d0}
We say that a mechanism $X$ meets $D_0$ if it satisfies all four of the
following.
\begin{itemize}
\item[$I_1$] (\emph{prime-distinguishing}) the asserted property depends on the
individual primes (the Euler product), not merely on prime counts or on the
functional equation and growth of $\zeta$;
\item[$I_2$] (\emph{independent input with a decisive theorem}) $X$ lies in a
class $\mathcal C$ carrying a general theorem $T$, established independently of
RH and holding for \emph{every} object of $\mathcal C$, such that $T$ applied
to $X$ yields the bound off the line as a \emph{corollary} ($I_{2a}$: such
$\mathcal C$ and $T$ exist; $I_{2b}$: the bound is a corollary of $T$, not a
property special to $X$);
\item[$I_3$] (\emph{individual-resolving}) the property constrains each single
zero, not only a sum, an integral, a moment or a density of zeros;
\item[$I_4$] (\emph{arithmetic-aware}) the property is not preserved when the
primes are replaced by another sequence with the same counting asymptotics.
\end{itemize}
\end{criterion}

The clauses are informal, and deciding whether a given programme meets them is
a matter of judgment. Clause $I_1$ is motivated by an example in which the
conclusion fails: the Davenport--Heilbronn function satisfies a functional
equation of the same type as $\zeta$, has no Euler product, and has zeros off
the critical line \cite[\S10.25]{bgTitchmarsh}. Clause $I_4$ expresses the same
requirement for mechanisms that see the primes only through their counting
function.

\begin{remark}[The role of $I_{2b}$]
Clause $I_{2b}$ is the one that separates the function-field case from the
rest. For the surface $C\times C$, with $C$ a curve over $\bbF_q$, the bound
$|\alpha_i|=\sqrt q$ for the Frobenius eigenvalues is a corollary of the Hodge
index theorem, which holds for \emph{all} smooth projective surfaces
\cite{Weil48}; the object lies in a class where a general theorem applies. For
$\zeta$, the analogous positivity (of Weil's quadratic functional in the
explicit formula \cite{Weil52}) is not a corollary of any general theorem
proved elsewhere; it is equivalent to RH.
\end{remark}

\begin{remark}[Assessment of known programmes]\label{prop:cal}
Table~\ref{tab:cal} lists the programmes we examined and, for each, a
clause of $D_0$ that in our assessment it does not meet (the one we regard
as decisive). The assessment is a
judgment, not a theorem. In this assessment the only mechanisms meeting all
four clauses are the proofs of RH over finite fields (Weil;
Stepanov--Bombieri; Deligne), and for every other programme clause $I_2$ is
not met or is open, or clause $I_3$ is not met.
\end{remark}

\begin{table}[ht]
\small
\begin{tabular}{>{\raggedright\arraybackslash}p{4.7cm}>{\raggedright\arraybackslash}p{1.5cm}>{\raggedright\arraybackslash}p{8.6cm}}
\hline
Programme & Clause not met & Reason \\
\hline
Weil's criterion \cite{Weil52,bgBombieri00} & $I_{2b}$ &
The positivity is equivalent to RH; no general theorem yields it.\\
Hilbert--P\'olya; Berry--Keating \cite{BerryKeating99} & $I_{2a}$ &
No self-adjoint operator is known whose spectrum is shown independently to be
the set of ordinates.\\
Connes' trace formula \cite{bgConnes99} & $I_{2b}$ &
The validity of the global trace formula is equivalent to RH.\\
Deninger's cohomological formalism \cite{Deninger98} & $I_2$ (open) &
The cohomology theory with the required properties is conjectural.\\
Connes--Consani \cite{ConnesConsani21} & $I_2$ (open) &
A geometric framework is under construction; no analogue of the Hodge index
theorem has been established in it.\\
de Branges spaces \cite{bgdeBranges} & $I_{2b}$ &
The positivity conditions that would imply RH are properties of the specific
space, not consequences of the general theory.\\
Nyman--Beurling; B\'aez-Duarte \cite{Beurling55,BaezDuarte03} & $I_{2b}$ &
The density statement is equivalent to RH.\\
Li's criterion \cite{Li97} & $I_{2b}$ &
The positivity of the sequence $\lambda_n$ is equivalent to RH.\\
de Bruijn--Newman constant \cite{deBruijn50,Newman76,RodgersTao20} & $I_{2b}$ &
RH is equivalent to $\Lambda_{\mathrm{dBN}}\le0$, while
$\Lambda_{\mathrm{dBN}}\ge0$ is known.\\
Jensen polynomials, Tur\'an inequalities \cite{Polya27,CNV86,GORZ19} & $I_{2b}$ &
The full hierarchy is equivalent to RH (see \S\ref{sec:gap}).\\
Pair correlation \cite{bgMont73} & $I_3$ &
A statement on the statistics of zeros.\\
Random matrix models, moments \cite{bgKS,KeatingSnaith00} & $I_3$ &
Statistical models and conjectures.\\
Multiplicative chaos \cite{SaksmanWebb20} & $I_3$ &
A statement on the statistical behaviour of $\zeta$ on the line.\\
Zero-density estimates \cite[Ch.~9]{bgTitchmarsh} & $I_3$ &
They bound the number of zeros off the line and do not exclude a single one.\\
Levinson's method \cite{Levinson74,Conrey89} & $I_3$ &
It gives a proportion of zeros on the line.\\
Curves over $\bbF_q$: Weil \cite{Weil48} & none &
Hodge index theorem on $C\times C$.\\
Curves over $\bbF_q$: Stepanov--Bombieri \cite{Stepanov69,Bombieri74} & none &
Riemann--Roch on $C$ and an auxiliary function.\\
Varieties over $\bbF_q$: Deligne \cite{Deligne74,Deligne80} & none &
Purity of the cohomology.\\
\hline
\end{tabular}
\medskip
\caption{Programmes and a clause of $D_0$ not met (the author's
assessment).}\label{tab:cal}
\end{table}

Three situations occur in Table~\ref{tab:cal}: all clauses are met (a known
general theorem $T$ applies: finite fields); clause $I_2$ is open (a class is
being built in which a general $T$ would apply, but $T$ is not established:
Deninger, Connes--Consani); a clause is not met (the required statement is
equivalent to RH for $\zeta$, or it concerns the zeros only in aggregate).

\begin{remark}[The available general theorems]
The general theorems we know of that have delivered a bound of this kind as a
corollary are positivity theorems of Hodge type: the Hodge index theorem for
surfaces, purity for varieties over finite fields \cite{Deligne74,Deligne80}, the
Hodge--Riemann relations (in K\"ahler geometry and, after
Adiprasito--Huh--Katz \cite{AHK18} and Br\"and\'en--Huh \cite{BrandenHuh20},
for matroids and Lorentzian polynomials); to these one may add the spectral
theorem for a self-adjoint operator whose spectrum is known independently, and
Stepanov's method \cite{Stepanov69,Bombieri74}. For $\zeta$ the difficulties
are of two kinds: no geometric object over $\operatorname{Spec}\bbZ$ is known
to which such a theorem applies (no surface, no cohomology, no operator
constructed independently of the zeros), and in the known positivity
formulations the positivity is itself equivalent to RH. The next two sections
illustrate the second and the first of these, respectively.
\end{remark}

\section{Finite orders and the full order}\label{sec:gap}

Let $\Xib(x)=\xi(\tfrac12+ix)$. Recall P\'olya's kernel
\[
\Phi(t)=\sum_{n\ge1}\bigl(2\pi^2 n^4 e^{9t/2}-3\pi n^2 e^{5t/2}\bigr)
e^{-\pi n^2 e^{2t}},
\]
which is even, positive, and satisfies
$\Xib(x)=2\int_{-\infty}^{\infty}\Phi(t)e^{ixt}\,dt$
\cite[\S10.1]{bgTitchmarsh}, \cite{CNV86}. Thus $\Xib$ is real and even, and RH
holds if and only if $\Xib$ belongs to the Laguerre--P\'olya class $\LP$ (real
entire functions that are locally uniform limits of real polynomials with only
real zeros). Put $m_{2k}=\int_{-\infty}^{\infty}t^{2k}\Phi(t)\,dt$ and
$b(k)=m_{2k}/(2k)!$, so that $\Xib(x)=2\sum_{k\ge0}(-1)^k b(k)x^{2k}$.

Following Schoenberg \cite{Schoenberg51}, a \emph{P\'olya frequency function}
is a measurable function $\Lambda$ on $\bbR$ with $0<\int\Lambda<\infty$ such
that $\det[\Lambda(x_i-y_j)]_{i,j=1}^N\ge0$ for all $N\ge1$ and all
$x_1<\dots<x_N$, $y_1<\dots<y_N$. A function is $\TP_k$ if these
determinants are nonnegative for all $N\le k$. Since $\xi(\tfrac12+y)$ is
positive for real $y$ and $1/\xi(\tfrac12+y)$ is holomorphic and rapidly
decreasing in the strip $|\Im y|<\gamma_1$, where $\gamma_1=14.13\ldots$ is
the ordinate of the first zero, the function
\[
\Lambda(x)=\frac1{2\pi}\int_{-\infty}^{\infty}\frac{e^{ixy}}{\xi(\tfrac12+y)}\,dy
\]
is well defined, real, even, continuous and integrable (shifting the contour
gives $\Lambda(x)=O(e^{-c|x|})$ for every $c<\gamma_1$), and
$\int_{-\infty}^{\infty}\Lambda(x)e^{-sx}\,dx=1/\Xib(s)$ for
$|\Re s|<\gamma_1$.

\begin{theorem}[Total-positivity form of RH]\label{thm:tp}
$\RH$ holds if and only if $\Lambda$ is a P\'olya frequency function.
\end{theorem}
\begin{proof}
By Schoenberg's theorem \cite{Schoenberg51}, the bilateral Laplace transform
of a P\'olya frequency function converges in an open vertical strip containing
the imaginary axis and equals $1/\Psi(s)$ there, with
\[
\Psi(s)=Ce^{-cs^2+\delta s}\prod_\nu(1+\delta_\nu s)e^{-\delta_\nu s},\qquad
C>0,\ c\ge0,\ \delta,\delta_\nu\in\bbR,\ 0<c+\sum_\nu\delta_\nu^2<\infty ;
\]
conversely, for every such $\Psi$ the function $1/\Psi$ is, in such a strip,
the Laplace transform of a P\'olya frequency function. If $\Lambda$ is a
P\'olya frequency function, then $\Xib=\Psi$ on the imaginary axis, hence
everywhere, and all zeros of $\Xib$ are real. Conversely, $\Xib$ is real,
even, of order $1$, with $\Xib(0)>0$ and $\sum|\gamma|^{-2}<\infty$, so that
$\Xib(s)=\Xib(0)\prod_{\Re\gamma>0}(1-s^2/\gamma^2)$; if all $\gamma$ are real
this is of the form $\Psi$ with $c=\delta=0$ and
$\{\delta_\nu\}=\{\pm1/\gamma\}$, so $1/\Xib$ is the Laplace transform of a
P\'olya frequency function, which coincides with $\Lambda$ by uniqueness of
the Fourier transform.
\end{proof}

This formulation of RH is discussed by Gr\"ochenig \cite{Groechenig20}; for
related results on multiply positive sequences see Katkova \cite{Katkova07}.
The kernel $\Phi$ itself is \emph{not} a P\'olya frequency function: it is
integrable and its Laplace transform on the imaginary axis is
$\tfrac12\Xib$, which has zeros, whereas $1/\Psi$ has none. It is $\TP_2$,
since for a translation kernel $\TP_2$ is log-concavity (see
Proposition~\ref{prop:orders}). Numerically, however, it fails $\TP_5$:
for instance $\det[\Phi(0.1(i-j))]_{i,j=1}^{5}\approx-1.39\cdot10^{-7}\,\Phi(0)^5$.

A second hierarchy is attached to the Taylor coefficients of $\Xib$. The
following collects the known unconditional results.

\begin{proposition}[Known unconditional orders]\label{prop:orders}
Let $c_k=k!\,m_{2k}/(2k)!=k!\,b(k)$ and
$J^{d,n}(X)=\sum_{j=0}^{d}\binom dj c_{n+j}X^j$. A real polynomial is called
hyperbolic if all its zeros are real. Then:
\begin{itemize}
\item[\rm(i)] $\Phi>0$ and $\Phi$ is log-concave on $\bbR$
{\rm(Csordas--Varga \cite{CV90})};
\item[\rm(ii)] $c_k^2\ge c_{k-1}c_{k+1}$ for all $k\ge1$;
equivalently $b(k)^2\ge\frac{k+1}{k}\,b(k-1)b(k+1)$; equivalently $J^{2,n}$ is
hyperbolic for all $n\ge0$ {\rm(the Tur\'an inequalities,
Csordas--Norfolk--Varga \cite{CNV86})};
\item[\rm(iii)] for each $d$, $J^{d,n}$ is hyperbolic for all $n\ge N(d)$, and
for all $n$ if $d\le 8$ {\rm(Griffin--Ono--Rolen--Zagier \cite{GORZ19})};
\item[\rm(iv)] $\RH$ holds if and only if $J^{d,n}$ is hyperbolic for all
$d\ge1$ and $n\ge0$ {\rm(P\'olya--Jensen \cite{Polya27})}.
\end{itemize}
\end{proposition}
\begin{proof}
For (i): $\Phi'(t)<0$ for $t>0$ \cite{CNV86} and $t\mapsto\log\Phi(\sqrt t)$ is
concave on $(0,\infty)$ \cite{CV90}. Writing $g(t)=\log\Phi(\sqrt t)$, one has
$(\log\Phi)''(u)=4u^2g''(u^2)+2g'(u^2)\le0$ for $u>0$, and $\Phi$ is even. The
equivalences in (ii) follow from $c_k=k!\,b(k)$ and from the discriminant
$4(c_{n+1}^2-c_nc_{n+2})$ of $J^{2,n}$. The remaining
statements are in the references; the normalisation of $c_k$ used there
differs from ours by a factor $c\,r^k$ with $c>0$ and $r\ne0$, which does not
affect hyperbolicity.
\end{proof}

Note that the Tur\'an inequalities in (ii) are stronger than log-concavity of
the sequence $b(k)$, and that (i) and (ii) are different statements: (i)
concerns the kernel, (ii) its moments.

\begin{remark}[Finite orders versus the full order]\label{thm:gap}
In the Jensen hierarchy the full condition (all $d$, all $n$) is equivalent to
RH, the conditions $d\le 8$ are known for all $n$, and for each $d$ the
condition is known for $n\ge N(d)$ (Proposition~\ref{prop:orders}). What is
missing is uniformity in $(d,n)$: no known theorem passes from the
fixed-degree statements to the full one. A similar pattern occurs in other
positivity formulations. Weil's functional \cite{Weil52} is known to be
positive on test functions supported in a sufficiently short interval
\cite{Yoshida92,ConnesConsani21}, and RH is equivalent to its positivity on
all test functions. For the function $\Lambda$ of Theorem~\ref{thm:tp}, RH is
equivalent to $\TP_k$ for all $k$; we are not aware of unconditional results
on $\TP_k$ for $\Lambda$. By contrast, not every positivity attached to
$\Phi$ is of this kind: the Hankel matrices $[m_{2(i+j)}]_{i,j=0}^{N}$ are
positive definite for every $N$, because
$\sum_{i,j}c_ic_jm_{2(i+j)}=\int(\sum_ic_it^{2i})^2\Phi(t)\,dt>0$, and this
holds whether or not RH does. We do not claim an order-by-order equivalence
between these hierarchies.
\end{remark}

\begin{remark}
In Proposition~\ref{prop:orders} the parameter $n$ is the shift of the Jensen
polynomial, not a height on the critical line: each degree $d$ is settled for
large shift, and RH requires all degrees at all shifts.
\end{remark}

\section{Periodicity of the zeros over finite fields, and its absence for
\texorpdfstring{$\zeta$}{zeta}}\label{sec:fin}

We recall why, for curves over finite fields, the positivity equivalent to RH
is a condition on finitely many numbers, and observe that the feature
responsible for this has no counterpart for $\zeta$.

\begin{proposition}[Geometry makes the positivity finite]\label{prop:finit}
Let $C/\bbF_q$ be a smooth projective geometrically irreducible curve of genus
$g\ge1$, with zeta function $Z(C,T)=P(T)/((1-T)(1-qT))$,
$P(T)=\prod_{i=1}^{2g}(1-\alpha_iT)$. RH for $C$ is a \emph{finite} condition:
$|\alpha_i|=\sqrt q$ for the $2g$ numbers $\alpha_i$, and it follows from the
Hodge index theorem on the lattice $\NS(C\times C)$ of finite rank. Writing
$\alpha_i=\sqrt q\,e^{i\theta_i}$ with $\theta_i\in\bbR$, the zeros of
$P(q^{-s})$, counted with multiplicity, are
\[
s=\tfrac12+i\frac{\theta_i+2\pi k}{\log q}\qquad(1\le i\le 2g,\ k\in\bbZ):
\]
infinitely many zeros, but only $2g$ modulo the period $2\pi i/\log q$; the
ordinates are periodic with period $2\pi/\log q$.
\end{proposition}
\begin{proof}
The Hodge index theorem on $C\times C$, applied to the classes of the graphs
of the powers of the Frobenius morphism, gives
$|\#C(\bbF_{q^r})-q^r-1|\le 2g\,q^{r/2}$ for all $r\ge1$, and this is
equivalent to $|\alpha_i|=\sqrt q$ for all $i$ \cite{Weil48}; see
\cite[V, Ex.~1.10 and App.~C, Ex.~5.7]{Hartshorne}. The zeros of $P(q^{-s})$
are the solutions of $q^{s}=\alpha_i$, which are the numbers displayed.
\end{proof}

In the function-field case the device that makes the positivity a finite
condition is this periodicity. We call a structure of this kind a
\emph{finitization}: a reduction of the positivity equivalent to RH to a
positivity statement on a module of finite rank.

\begin{proposition}[No Frobenius-type periodicity]\label{thm:LI}
There are no $\tau>0$ and finite set $\{\beta_1,\dots,\beta_m\}\subset\bbR$
such that every ordinate $\gamma>0$ of a zero of $\zeta$ on the critical line
lies in $\{\beta_j+k\tau:1\le j\le m,\ k\in\bbZ\}$.
\end{proposition}
\begin{proof}
The set $\{\beta_j+k\tau\}$ has at most $m(T/\tau+1)$ points in $(0,T]$. On
the other hand a positive proportion of the zeros of $\zeta$ are simple and
lie on the critical line \cite{HeathBrown79} (see also
\cite{Selberg42,Levinson74,Conrey89}), and the number of zeros with ordinate
in $(0,T]$ is asymptotic to $\frac{T}{2\pi}\log T$
\cite[Thm.~9.4]{bgTitchmarsh}. Distinct simple zeros on the line have distinct
ordinates, so there are at least $cT\log T$ distinct ordinates in $(0,T]$ for
some $c>0$ and all large $T$. This exceeds $m(T/\tau+1)$ for large $T$.
\end{proof}

\begin{remark}[The formal limit $q\to1$]\label{cor:degen}
The period $2\pi/\log q$ of the ordinates in Proposition~\ref{prop:finit}
tends to infinity as $q\to1^+$. In the heuristic picture of $\operatorname{Spec}\bbZ$ as a curve
over $\bbF_1$, the fundamental domain containing finitely many zeros therefore
becomes the whole line, in agreement with Proposition~\ref{thm:LI}.
\end{remark}

\begin{remark}[Linear independence]
The Linear Independence conjecture \cite{Ingham42} asserts that the positive
ordinates of the zeros of $\zeta$ are linearly independent over $\bbQ$. Under
it, the conclusion of Proposition~\ref{thm:LI} follows from three ordinates
only. Indeed there are infinitely many distinct ordinates and finitely many
classes $\beta_j+\tau\bbZ$, so three distinct ordinates
$\gamma_{n_1},\gamma_{n_2},\gamma_{n_3}$ lie in one class; writing
$\gamma_{n_1}-\gamma_{n_2}=k\tau$ and $\gamma_{n_1}-\gamma_{n_3}=k'\tau$ with
$k,k'\in\bbZ\setminus\{0\}$ gives
$(k'-k)\gamma_{n_1}-k'\gamma_{n_2}+k\gamma_{n_3}=0$, a nontrivial rational
relation. Proposition~\ref{thm:LI} shows that the hypothesis is not needed:
the obstruction to a periodicity is the density of the zeros, which comes
from the archimedean factor. Periodicity and independence are in any case not
opposed: for an ordinary elliptic curve over $\bbF_q$ the zeros are periodic,
while $\theta_1/\pi$ is irrational.
\end{remark}

\section{Conclusion}
In the positivity formulations of RH considered in \S\ref{sec:gap}, finitely
many conditions, or conditions of each fixed order in a restricted range, are
known unconditionally, while the full family of conditions is equivalent to
RH (Theorem~\ref{thm:tp}, Proposition~\ref{prop:orders},
Remark~\ref{thm:gap}). Over finite fields the corresponding positivity is a
condition on finitely many numbers, because the zeros are periodic
(Proposition~\ref{prop:finit}), and it is supplied by a general theorem of
geometry. For $\zeta$ no such periodicity is possible
(Proposition~\ref{thm:LI}), and the period of the function-field case tends
to infinity in the formal limit $q\to1$ (Remark~\ref{cor:degen}). The
checklist of \S\ref{sec:disc} summarises this comparison: in our assessment,
the programmes for $\zeta$ either rest on a statement equivalent to RH, or
concern the zeros in aggregate, or await a general positivity theorem in a
class containing $\zeta$. The theorems quoted here are known; the note
contributes only their arrangement and the checklist, and it is not progress
towards RH.

\begin{thebibliography}{99}
\bibitem{AHK18} K.~Adiprasito, J.~Huh and E.~Katz, \emph{Hodge theory for combinatorial geometries}, Ann.\ of Math.\ (2) \textbf{188} (2018), 381--452.
\bibitem{BaezDuarte03} L.~B\'aez-Duarte, \emph{A strengthening of the Nyman--Beurling criterion for the Riemann hypothesis}, Atti Accad.\ Naz.\ Lincei Rend.\ Lincei Mat.\ Appl.\ \textbf{14} (2003), 5--11.
\bibitem{BerryKeating99} M.~V.~Berry and J.~P.~Keating, \emph{The Riemann zeros and eigenvalue asymptotics}, SIAM Rev.\ \textbf{41} (1999), 236--266.
\bibitem{Beurling55} A.~Beurling, \emph{A closure problem related to the Riemann zeta-function}, Proc.\ Nat.\ Acad.\ Sci.\ U.S.A.\ \textbf{41} (1955), 312--314.
\bibitem{bgBombieri00} E.~Bombieri, \emph{Problems of the Millennium: The Riemann Hypothesis}, Clay Mathematics Institute, 2000.
\bibitem{Bombieri74} E.~Bombieri, \emph{Counting points on curves over finite fields (d'apr\`es S.~A.~Stepanov)}, S\'eminaire Bourbaki 1972/73, Exp.\ 430, Lecture Notes in Math.\ \textbf{383}, Springer, 1974, 234--241.
\bibitem{BrandenHuh20} P.~Br\"and\'en and J.~Huh, \emph{Lorentzian polynomials}, Ann.\ of Math.\ (2) \textbf{192} (2020), 821--891.
\bibitem{deBruijn50} N.~G.~de Bruijn, \emph{The roots of trigonometric integrals}, Duke Math.\ J.\ \textbf{17} (1950), 197--226.
\bibitem{bgdeBranges} L.~de Branges, \emph{Hilbert Spaces of Entire Functions}, Prentice-Hall, 1968.
\bibitem{bgConnes99} A.~Connes, \emph{Trace formula in noncommutative geometry and the zeros of the Riemann zeta function}, Selecta Math.\ (N.S.) \textbf{5} (1999), 29--106.
\bibitem{ConnesConsani21} A.~Connes and C.~Consani, \emph{Weil positivity and trace formula, the archimedean place}, Selecta Math.\ (N.S.) \textbf{27} (2021), Paper No.~77.
\bibitem{bgConrey03} J.~B.~Conrey, \emph{The Riemann Hypothesis}, Notices Amer.\ Math.\ Soc.\ \textbf{50} (2003), 341--353.
\bibitem{Conrey89} J.~B.~Conrey, \emph{More than two fifths of the zeros of the Riemann zeta function are on the critical line}, J.\ Reine Angew.\ Math.\ \textbf{399} (1989), 1--26.
\bibitem{CNV86} G.~Csordas, T.~S.~Norfolk and R.~S.~Varga, \emph{The Riemann hypothesis and the Tur\'an inequalities}, Trans.\ Amer.\ Math.\ Soc.\ \textbf{296} (1986), 521--541.
\bibitem{CV90} G.~Csordas and R.~S.~Varga, \emph{Necessary and sufficient conditions and the Riemann hypothesis}, Adv.\ in Appl.\ Math.\ \textbf{11} (1990), 328--357.
\bibitem{Deligne74} P.~Deligne, \emph{La conjecture de Weil.\ I}, Publ.\ Math.\ Inst.\ Hautes \'Etudes Sci.\ \textbf{43} (1974), 273--307.
\bibitem{Deligne80} P.~Deligne, \emph{La conjecture de Weil.\ II}, Publ.\ Math.\ Inst.\ Hautes \'Etudes Sci.\ \textbf{52} (1980), 137--252.
\bibitem{Deninger98} C.~Deninger, \emph{Some analogies between number theory and dynamical systems on foliated spaces}, Doc.\ Math., Extra Vol.\ ICM Berlin 1998, vol.~I, 163--186.
\bibitem{bgEdwards} H.~M.~Edwards, \emph{Riemann's Zeta Function}, Academic Press, 1974.
\bibitem{GORZ19} M.~Griffin, K.~Ono, L.~Rolen and D.~Zagier, \emph{Jensen polynomials for the Riemann zeta function and other sequences}, Proc.\ Natl.\ Acad.\ Sci.\ USA \textbf{116} (2019), 11103--11110.
\bibitem{Groechenig20} K.~Gr\"ochenig, \emph{Schoenberg's theory of totally positive functions and the Riemann zeta function}, arXiv:2007.12889; in: Sampling, Approximation, and Signal Analysis, Birkh\"auser, 2024, 193--210.
\bibitem{Hartshorne} R.~Hartshorne, \emph{Algebraic Geometry}, Graduate Texts in Math.\ \textbf{52}, Springer, 1977.
\bibitem{HeathBrown79} D.~R.~Heath-Brown, \emph{Simple zeros of the Riemann zeta-function on the critical line}, Bull.\ London Math.\ Soc.\ \textbf{11} (1979), 17--18.
\bibitem{Ingham42} A.~E.~Ingham, \emph{On two conjectures in the theory of numbers}, Amer.\ J.\ Math.\ \textbf{64} (1942), 313--319.
\bibitem{Katkova07} O.~M.~Katkova, \emph{Multiple positivity and the Riemann zeta-function}, Comput.\ Methods Funct.\ Theory \textbf{7} (2007), 13--31; arXiv:math/0505174.
\bibitem{bgKS} N.~M.~Katz and P.~Sarnak, \emph{Zeroes of zeta functions and symmetry}, Bull.\ Amer.\ Math.\ Soc.\ \textbf{36} (1999), 1--26.
\bibitem{KeatingSnaith00} J.~P.~Keating and N.~C.~Snaith, \emph{Random matrix theory and $\zeta(1/2+it)$}, Comm.\ Math.\ Phys.\ \textbf{214} (2000), 57--89.
\bibitem{Levinson74} N.~Levinson, \emph{More than one third of zeros of Riemann's zeta-function are on $\sigma=1/2$}, Adv.\ Math.\ \textbf{13} (1974), 383--436.
\bibitem{Li97} X.-J.~Li, \emph{The positivity of a sequence of numbers and the Riemann hypothesis}, J.\ Number Theory \textbf{65} (1997), 325--333.
\bibitem{bgMont73} H.~L.~Montgomery, \emph{The pair correlation of zeros of the zeta function}, Proc.\ Sympos.\ Pure Math.\ \textbf{24} (1973), 181--193.
\bibitem{Newman76} C.~M.~Newman, \emph{Fourier transforms with only real zeros}, Proc.\ Amer.\ Math.\ Soc.\ \textbf{61} (1976), 245--251.
\bibitem{Polya27} G.~P\'olya, \emph{\"Uber die algebraisch-funktionentheoretischen Untersuchungen von J.~L.~W.~V.~Jensen}, Kgl.\ Danske Vid.\ Sel.\ Math.-Fys.\ Medd.\ \textbf{7} (1927), no.~17, 3--33.
\bibitem{RodgersTao20} B.~Rodgers and T.~Tao, \emph{The de Bruijn--Newman constant is non-negative}, Forum Math.\ Pi \textbf{8} (2020), e6.
\bibitem{SaksmanWebb20} E.~Saksman and C.~Webb, \emph{The Riemann zeta function and Gaussian multiplicative chaos: statistics on the critical line}, Ann.\ Probab.\ \textbf{48} (2020), 2680--2754.
\bibitem{Schoenberg51} I.~J.~Schoenberg, \emph{On P\'olya frequency functions.\ I.\ The totally positive functions and their Laplace transforms}, J.\ Analyse Math.\ \textbf{1} (1951), 331--374.
\bibitem{Selberg42} A.~Selberg, \emph{On the zeros of Riemann's zeta-function}, Skr.\ Norske Vid.-Akad.\ Oslo I (1942), no.~10.
\bibitem{Stepanov69} S.~A.~Stepanov, \emph{On the number of points of a hyperelliptic curve over a finite prime field}, Izv.\ Akad.\ Nauk SSSR Ser.\ Mat.\ \textbf{33} (1969), 1171--1181.
\bibitem{bgTitchmarsh} E.~C.~Titchmarsh, \emph{The Theory of the Riemann Zeta-Function}, 2nd ed., rev.\ D.~R.~Heath-Brown, Oxford Univ.\ Press, 1986.
\bibitem{Weil48} A.~Weil, \emph{Sur les courbes alg\'ebriques et les vari\'et\'es qui s'en d\'eduisent}, Hermann, Paris, 1948.
\bibitem{Weil52} A.~Weil, \emph{Sur les ``formules explicites'' de la th\'eorie des nombres premiers}, Comm.\ S\'em.\ Math.\ Univ.\ Lund, Tome Suppl\'ementaire (1952), 252--265.
\bibitem{Yoshida92} H.~Yoshida, \emph{On Hermitian forms attached to zeta functions}, in: Zeta Functions in Geometry, Adv.\ Stud.\ Pure Math.\ \textbf{21}, Kinokuniya, 1992, 281--325.
\end{thebibliography}

\end{document}
